\hypertarget{appendix-1652d27a8e90}{}
\section*{Half adder}
\chapter{Half adder}\label{appendix-1652d27a8e90-half-adder}

\textbf{Curriculum depth:} Introduced · \textbf{First seen:} C002 · \textbf{Latest development:} C002

\section{Definition and mechanism}\label{appendix-1652d27a8e90-definition-and-mechanism}

Adds two bits with \(S=A\oplus B\) and \(C=A\land B\), but accepts no incoming carry.

A half adder adds two bits with sum S=A XOR B and carry C=A AND B. Inputs 1,1 produce binary 10: low sum 0, carry 1. It cannot consume a carry from a previous position, which motivates the full adder.

\section{Prerequisites}\label{appendix-1652d27a8e90-prerequisites}

\begin{itemize}
\tightlist
\item
  \hyperref[appendix-d49161831368]{Bitwise XOR} --- XOR returns 1 when inputs differ.
\item
  \hyperref[appendix-68b52453737c]{Bitwise AND} --- Per lane: 0\&0=0, 0\&1=0, 1\&0=0, 1\&1=1.
\end{itemize}

\section{Used by}\label{appendix-1652d27a8e90-used-by}

\begin{itemize}
\tightlist
\item
  \hyperref[appendix-d1a5b1e520c2]{Full adder} --- uses half adder as part of its representation or reasoning.
\end{itemize}

\section{Recall and next depth}\label{appendix-1652d27a8e90-recall-and-next-depth}

Explain half adder without the label, work the example above, then state the assumption whose removal would change the result.

\textbf{Pending:} Apply and reconstruct this concept independently.

\section{Source and navigation}\label{appendix-1652d27a8e90-source-and-navigation}

\href{Continuous\%20Notes/Computing\%20Foundations\%20-\%20From\%20Bits\%20to\%20C++\%20and\%20Systems\%20APIs\#Arithmetic\%20becomes\%20a\%20policy}{C002 teaching chapter} · \hyperref[appendix-fd090d982a15]{Part definitions and mechanisms} · \href{../02-Foundations-II/main.pdf\#nameddest=appendix-4ad47bdd269a}{Dependency Map (02-Foundations-II, appendix-4ad47bdd269a)}

\section{From truth table to gate delay}\label{appendix-1652d27a8e90-from-truth-table-to-gate-delay}

A half adder adds two one-bit inputs \(a,b\). The sum bit is XOR and the carry bit is AND:

\[
s=a\oplus b,\qquad c=a\land b.
\]

The truth table is not just a memorization device; it proves why one output bit cannot encode all four input sums. Input \(1+1\) produces binary \(10\), requiring a second bit for carry.

A half adder has no carry-in, so it cannot by itself build every bit position of a multi-bit integer addition. A full adder incorporates carry-in and can be composed into ripple-carry or faster prefix structures. The important systems connection is dependency depth: in a ripple chain, each bit may wait for the previous carry, whereas parallel-prefix adders reorganize generate/propagate information to reduce critical path.

A reconstruction exercise should derive XOR/AND from the four rows, compose two half adders into a full adder, then explain why gate count and propagation depth are different optimization targets.
